\documentclass[a4paper, 11pt]{article}
\usepackage[catalan]{babel}
\usepackage{fontspec}
\setmainfont{Noto Sans}
\usepackage[table,xcdraw]{xcolor}
\pagestyle{fancy}
\usepackage{fancyhdr}
\usepackage{enumitem}
\usepackage{geometry}
\geometry{margin=2.5cm}

\lhead{\textbf{2n ESO B MAT \textcolor{red}{(PI)}}}
\chead{\textbf{Nombres i Fraccions}}
\rhead{23/03/2026}
\rfoot{\thepage}

\begin{document}

\begin{center}
    {\huge \textbf{EXAMEN ADAPTAT (PI)}} \\
    \vspace{0.5cm}
    \textbf{Alumne/a: IU MARTÍNEZ} \\
    \vspace{0.3cm}
    \textit{Pots fer servir la calculadora i el formulari de fraccions.}
\end{center}

\vspace{1cm}

\section*{PART A — Qüestions (4 punts)}

\begin{enumerate}
    \item Calcula: $-5 + 6 \cdot 2$
    \begin{enumerate}[label=\Alph*)] \item 7 \item 2 \item -11 \item 1 \end{enumerate}

    \item Si a la nevera estem a $-4^\circ\text{C}$ i pugem $11^\circ\text{C}$, a quina temperatura estarem?
    \begin{enumerate}[label=\Alph*)] \item $7^\circ\text{C}$ \item $15^\circ\text{C}$ \item $-7^\circ\text{C}$ \item $0^\circ\text{C}$ \end{enumerate}

    \item Quina fracció és equivalent a $\frac{1}{2}$?
    \begin{enumerate}[label=\Alph*)] \item $\frac{2}{4}$ \item $\frac{2}{3}$ \item $\frac{1}{4}$ \item $\frac{3}{2}$ \end{enumerate}

    \item Calcula $\frac{3}{4} + \frac{1}{4}$:
    \begin{enumerate}[label=\Alph*)] \item $\frac{4}{8}$ \item $\frac{4}{4} = 1$ \item $\frac{2}{4}$ \item $\frac{3}{16}$ \end{enumerate}
\end{enumerate}

\newpage

\section*{PART B — Problemes (6 punts)}

\noindent\textbf{Problema 1 (3 punts): L'ascensor} \\
Estem a la planta $-2$. L'ascensor puja 5 plantes i després en baixa 8. \\
\textbf{a)} A quina planta arribem al final? \\
\vspace{4cm}

\noindent\textbf{Problema 2 (3 punts): El dipòsit} \\
Tenim un dipòsit de 60 litres. N'hem consumit $\frac{1}{4}$ part. \\
\textbf{a)} Quants litres hem gastat? \\
\vspace{4cm}
\textbf{b)} Quants litres queden al dipòsit? \\
\vspace{4cm}

\end{document}
